\chapter*{Nota do Autor}
Essa história sempre viveu no fundo da minha mente, e eu sabia que um dia precisaria colocá-la no papel. Desde pequeno, quem me conhece sabe que nutro a esperança de poder voar com minhas próprias forças, como um super-herói. Não era apenas a vontade de sair do chão. Era também a fantasia de escapar de tudo aquilo que parecia pesado demais para enfrentar andando.

Escrever este livro até o fim foi, para mim, uma espécie de processo terapêutico. Tentei construir personagens vivos, com desejos, medos, contradições e anseios próprios. Depois, ao reler a história, notei que havia um pouco de mim em cada um deles: no Garoto que acredita poder resolver tudo com superpoderes; em quem foge; em quem tenta controlar; em quem transforma a própria dor numa certeza sobre a vida dos outros.

Talvez o Garoto demore a perceber que os grandes problemas da vida fazem parte da nossa jornada. Lidar com cada momento usando a força que já temos pode exigir muito mais coragem do que sair voando de uma situação ou ficar invisível para que ninguém nos veja. Nem sempre existe uma habilidade especial capaz de apagar o medo, a ansiedade ou a angústia. Às vezes, o gesto mais poderoso é continuar presente.

Durante o tempo que levei para escrever, reler e revisar este livro, pude refletir sobre por que eu queria resolver as coisas com poderes, em vez de encará-las de frente. Hoje me sinto um pouco mais preparado para lidar com minhas ansiedades e angústias usando as ferramentas que tenho, mesmo quando elas parecem pequenas. Essa mudança não resolve tudo, mas me ajuda a reconhecer o que antes eu preferia evitar.



\begin{tikzpicture}[remember picture, overlay]
    \node[anchor=south west, align=left] at (current page text area.south west) {Apesar de, no fundo, eu ainda acreditar que um dia poderei voar.};
  \end{tikzpicture}
  
  
  
  


